\documentclass[10pt,a4paper]{amsart}
\usepackage[margin=19mm]{geometry}
\usepackage{amsmath,amssymb,mathtools}
\usepackage[scr=rsfs,cal=euler]{mathalfa}
\usepackage{xfrac}
\usepackage[unicode,hidelinks,bookmarks=false]{hyperref}
\newcommand{\ie}{i.e.\ }
\newcommand{\cf}{cf.\ }
\newcommand{\resp}{respectively\ }
\newcommand{\Subsection}{\subsection}
\providecommand{\nobreakdash}{\nobreak\hbox{-}}
\setlength{\emergencystretch}{2em}
\multlinegap0pt
\usepackage{bm}
\usepackage{bbm}
\usepackage{dsfont}
\usepackage{xspace}
\usepackage{cancel}
\usepackage{mathtools}
\usepackage{amsthm}
\makeatletter
\@ifundefined{theorem}{\newtheorem{theorem}{Theorem}}{}
\@ifundefined{assumption}{\newtheorem{assumption}{Assumption}}{}
\@ifundefined{lemma}{\newtheorem{lemma}[theorem]{Lemma}}{}
\makeatother
\newcommand{\E}{\mathbb{E}}
\newcommand{\R}{\mathbb{R}}
\renewcommand{\P}{\mathbb{P}} 
\title[Dropout and Random Gradient Masking Are Asymptotically Equiv]{Dropout and Random Gradient Masking Are Asymptotically Equivalent in Large ResNets}
\author{Javier Maass, L\'ena\"ic Chizat}
\date{Source: arXiv:2607.16761v1 (CC BY 4.0)}
\begin{document}
\maketitle

\noindent\emph{Source and licence.} Dropout and Random Gradient Masking Are Asymptotically Equivalent in Large ResNets. Version arXiv:2607.16761v1 (CC BY 4.0), distributed under the Creative Commons Attribution 4.0 International licence, as recorded in the arXiv licence metadata for this exact version and shown on the abstract page https://arxiv.org/abs/2607.16761v1. Full rights notice: licenses/arxiv-2607.16761-CC-BY-4.0.txt.\par\smallskip

\noindent\textit{Editorial scope.} Counted unit: the Lemma whose source label is \texttt{lem:propagation\_of\_subgaussianity} in the article (source0.tex lines 741--746, source label \texttt{lem:propagation\_of\_subgaussianity}), delivered together with the article's own complete original proof of it (source0.tex lines 748--764, one \texttt{proof} environment of 17 lines). No mathematics is written, repaired or re-derived: statement and proof are the article's own text line for line. Required context retained uncounted: ass:regularity\_assumptions (lines 353--360); eq:limit\_update\_rule\_ram (lines 458--460); lem:propagation\_of\_regularity (lines 687--695); lem:subgaussian\_via\_lipschitz (lines 722--726); lem:subgaussian\_and\_bounded\_independent (lines 729--732). Every definition, display and label of those lines is retained unchanged. Omitted from this package and not redistributed: 1--352, 361--457, 461--686, 696--721, 727--728, 733--740, 747--747, 765--1275. Nothing inside a retained excerpt is omitted or altered: every retained body is delivered exactly as the article's lines are written, so no omission receipt of any kind accompanies this package. Citations: the retained excerpts carry no citation command at all, so no bibliography entry of the article is transcribed and no reference is redistributed. Reference adaptation: none was needed and none was made; every reference inside the delivered unit resolves inside it, so no unresolved reference or citation is produced. Printed slips kept literal: no printed word, symbol, hypothesis or number of the counted statement or of its proof was altered. Layout and class: the article's own class is \texttt{article} with packages including geometry, inputenc, fontenc, babel, enumerate, enumitem, subcaption, xcolor, pdfpages, ulem, titling, hyperref, csquotes, tocbibind; the wrapper is the lane's pinned \texttt{amsart} editorial wrapper at 10pt on A4, which restates the theorem-like environments the retained text needs and adds the article's own macro definitions that the retained text needs (\texttt{\textbackslash E}, \texttt{\textbackslash P}, \texttt{\textbackslash R}), with extra wrapper packages (bm, bbm, dsfont, xspace, cancel, mathtools). The delivered numbering is the unit's own. Assets: no retained span contains a figure environment, a graphics-inclusion command, a PDF-page inclusion, a table or a TikZ picture; no image file is copied into this package, the package declares no asset dependency and no image file is redistributed with it. Licence: version arXiv:2607.16761v1 of this article is distributed under the Creative Commons Attribution 4.0 International licence, as recorded in the arXiv licence metadata for this exact version and shown on the abstract page https://arxiv.org/abs/2607.16761v1. A full rights notice is delivered at licenses/arxiv-2607.16761-CC-BY-4.0.txt. Characters: every retained excerpt is pure ASCII, so the delivered engine, XeLaTeX with its default Latin Modern text font, reports no missing character.
\begin{assumption}[Regularity assumptions]\label{ass:regularity_assumptions}
    There exists $B > 0$ such that:
\begin{enumerate}
    \item $\phi$ is $B$-Lipschitz, differentiable, its differential $D\phi$ is $B$-Lipschitz and $\|\phi(0,0)\|_2 \le B$;
    \item For some fixed dataset ${(x_i, y_i)}_{i=1}^n$ with $n\in \mathbb{N}$, $\pi = \frac{1}{n}\sum_{i=1}^n \delta_{(x_i, y_i)}$ , and the inputs satisfy $\max_i \|x_i\|_2 \le B$. For notational convenience, we will write \(\hat{h}_{k,i}^\ell := \hat{h}_{\theta_k}^{\zeta_k}(\ell, x_i)\), \(\tilde{h}_{k,i}^\ell := \hat{h}_{\tilde{\theta}_k}^{0}(\ell, x_i)\), and similarly for \(\hat{b}_{k,i}^\ell\), \(\tilde{b}_{k,i}^\ell\).
    \item The loss is differentiable in its first argument, such that $y\mapsto \nabla_1 \mathrm{loss}(y, y_i)$ is $B$-Lipschitz and $\|\nabla_1 \mathrm{loss}(0, y_i)\|_2 \le B$ $\forall i \in [n]$;
\end{enumerate}
\end{assumption}

\begin{equation}\label{eq:limit_update_rule_ram}
  Z_{k+1}(s) = Z_k(s) - \tau \E_{\pi}[\tilde{\mathfrak{u}}(Z_k(s), h_{\mu_k}(s, x), b_{\mu_k}(s, x, \bar{w}_{x,y}))](1+\zeta_k).
\end{equation}

\begin{lemma}[Propagation of regularity]\label{lem:propagation_of_regularity}
Let Assumption~\ref{ass:regularity_assumptions} hold with $B > 0$ and consider the limit dynamics~\eqref{eq:limit_update_rule_ram} with $\mu_0 \in \mathcal{P}(\R^p)$ with first moment bounded by $B$ {and $\nu_\zeta \in \mathcal{P}(\R^D)$ bounded by $B_\zeta$}. Then, there exists $c$ that only depends on $B$ {and $B_\zeta$} such that for all $k \geq 0$, $i \in [1:n]$,

\begin{enumerate}
    \item $Z_k$, $h_{k,i}$ and $b_{k,i}$ are uniquely well-defined;
    \item the functions $s \mapsto h_{k,i}(s)$ and $s \mapsto b_{k,i}(s)$ are $c$-Lipschitz and bounded in norm by $c$;
    \item the function $s \mapsto Z_k(s)$ is $(e^{ck\tau} - 1)$-Lipschitz (surely).
\end{enumerate}
\end{lemma}

\begin{lemma}\label{lem:subgaussian_via_lipschitz}
    If $X$ is a subgaussian vector in $\R^p$ and $f:\R^p\to \R^D$ is a $L$-Lipschitz function, then
    \[\|f(X)-\E[f(X)]\|_{\psi_2} \leq cL\sqrt{p}\|X - \E[X]\|_{\psi_2}, \]
    for an absolute constant $c>0$. Equivalently, \(\|f(X)\|_{vp} \leq cL\sqrt{p}\|X\|_{vp}\), where $\|\cdot\|_{vp}$ is the variance-proxy seminorm.
\end{lemma}

\begin{lemma}\label{lem:subgaussian_and_bounded_independent}
    If $X$ is a subgaussian vector in $\R^p$ and $Y$ is a random vector in $\R^{\tilde{p}}$ bounded by $B>0$ and independent of $X$. Let $f:\R^p\times\R^{\tilde{p}}\to \R^D$ be a (measurable) function. Then,
    \(\|f(X, Y)\|_{\psi_2} \leq \sup_{\|y\|\leq B}\|f(X,y)\|_{\psi_2}\). In terms of the variance-proxy\footnote{{Note that in general it is false that $\|f(X,Y)\|_{vp} \leq \sup_y \|f(X,y)\|_{vp}$ (unless $f(X,Y)$ is centered), as can be seen from choosing $f(x,y) = y$.}}, we have \[\|f(X, Y)\|_{vp}^2 \leq \sup_{\|y\|\leq B}\|f(X,y)\|_{vp}^2 + \frac{1}{4}\sup_{\|y\|, \|y'\|\leq B}\|\E[f(X,y)]-\E[f(X,y')]\|_{2}^2.\]
\end{lemma}

\begin{lemma}[Propagation of subgaussian tails]\label{lem:propagation_of_subgaussianity}
Let Assumption~\ref{ass:regularity_assumptions} hold with $B > 0$ and consider the limit dynamics~\eqref{eq:limit_update_rule_ram} with $\mu_0 \in \mathcal{P}(\R^p)$ and $\nu_\zeta \in \P(\R^D)$. Assume moreover that $\mu_0$ is subgaussian {and $\nu_\zeta$ bounded by $B_\zeta$}. Then there exists $c > 0$ that only depends on $B$, $B_\zeta$ and $D$ such that $\forall k \geq 0$ and $s \in [0,1]$,
\[
    \|Z_k(s)\|_{vp} \le e^{ck\tau}({1+} \|Z_0\|_{vp}).
\]
\end{lemma}

\begin{proof}
By definition, $Z_0$ is subgaussian. Recursively, we use the seminorm properties of the variance-proxy to see that
\begin{multline*}
    \|Z_{k+1}(s)\|_{vp} \le \|Z_k(s)\|_{vp} + \tau \max_{i \in [1:n]} \|\tilde{\mathfrak{u}}(Z_k(s), h_{k,i}(s), b_{k,i}(s){;\zeta_k})\|_{vp}\\
    \le \|Z_k(s)\|_{vp} + \tau \max_{i \in [1:n]}\sup_{\|\xi\|_{\infty}\leq B_\zeta} \|\tilde{\mathfrak{u}}(Z_k(s), h_{k,i}(s), b_{k,i}(s){;\xi})\|_{vp} \\{+ c\sup_{\|\xi\|_{\infty}, \|\xi'\|_{\infty}\leq B_\zeta} \|\E[\tilde{\mathfrak{u}}(Z_k(s), h_{k,i}(s), b_{k,i}(s){;\xi}) - \tilde{\mathfrak{u}}(Z_k(s), h_{k,i}(s), b_{k,i}(s){;\xi'})]\|_{2},}
\end{multline*}
where in the last inequality we used Lemma~\ref{lem:subgaussian_and_bounded_independent} since $Z_k$ and $\zeta_k$ are independent.
{Whenever $b$ is restricted to a ball of radius $R$ (which it is by Lemma~\ref{lem:propagation_of_regularity}) and $\|\xi\|_\infty \leq B_\zeta$, the function $z \mapsto \tilde{\mathfrak{u}}(z,h,b;\xi)$ is $BR(1+B_\zeta)$-Lipschitz, so we can apply Lemma~\ref{lem:subgaussian_via_lipschitz}} to get
\[{\max_{i \in [1:n]}\sup_{\|\xi\|_{\infty}\leq B_\zeta} \|\tilde{\mathfrak{u}}(Z_k(s), h_{k,i}(s), b_{k,i}(s){;\xi})\|_{vp} \leq cBR(1+B_\zeta)\|Z_k\|_{vp}}.\]
Similarly, since $\xi \mapsto \tilde{\mathfrak{u}}(z,h,b;\xi)$ is $BR$-Lipschitz from $\Vert\cdot\Vert_\infty$ to $\Vert\cdot\Vert_2$, we have
\begin{multline*}
    \sup_{\|\xi\|_{\infty}, \|\xi'\|_{\infty}\leq B_\zeta} \|\E[\tilde{\mathfrak{u}}(Z_k(s), h_{k,i}(s), b_{k,i}(s){;\xi}) - \tilde{\mathfrak{u}}(Z_k(s), h_{k,i}(s), b_{k,i}(s){;\xi'})]\|_{2} \\\leq BR\sup_{\|\xi\|_{\infty}, \|\xi'\|_{\infty}\leq B_\zeta}\|\xi - \xi'\|_\infty \leq 2BRB_\zeta.
\end{multline*}
All in all, absorbing $B$, $B_\zeta$ and $R$ into the constant $c >0$, this yields
\[{\|Z_{k+1}(s)\|_{vp} \le (1+\tau c)\|Z_k(s)\|_{vp} + \tau c.}\]
{By the discrete Gronwall lemma, we conclude that there exists $c$ that only depends on $B$, $B_\zeta$ and $D$ such that for $k \ge 0$ it holds $\|Z_k(s)\|_{vp} \leq e^{c k \tau} (1+\|Z_0\|_{vp})$}.
\end{proof}

\end{document}
